\documentclass[11pt,a4paper]{article}

% ---- Packages ----------------------------------------------------------------
\usepackage{amsmath,amssymb}
\usepackage{physics}
\usepackage{hyperref}
\usepackage[margin=2.5cm]{geometry}
\usepackage{graphicx}
\usepackage{xcolor}
\usepackage{booktabs}
\usepackage{enumitem}
\usepackage{mathtools}
\usepackage{slashed}
\usepackage{tikz}

% ---- Theorem-like environments -----------------------------------------------
\newtheorem{remark}{Remark}
\newtheorem{definition}{Definition}

% ---- Custom commands ---------------------------------------------------------
\newcommand{\Dslash}{\not\hspace{-1pt}D}
\newcommand{\overleftarrowD}{\overleftarrow{D}}
\newcommand{\Trc}{\operatorname{Tr}_c}
\newcommand{\Trsc}{\operatorname{Tr}_{sc}}
\newcommand{\sign}{\operatorname{sign}}
\newcommand{\znum}{\mathbb{Z}_n}
\newcommand{\cc}{\mathrm{c.c.}}

\newcommand{\code}[1]{\texttt{#1}}

% ---- Title -------------------------------------------------------------------
\title{
  \textbf{Physics Note: Pion Energy-Momentum Tensor Measurement}\\[4pt]
  \large PyQUDA Meson EMT Workflow -- Connected 3pt, Quark 1pt, and Gluon 1pt
}
\author{
  PyQUDA Collaboration
}
\date{\today}

\begin{document}
\maketitle
\tableofcontents
\newpage

% ==============================================================================
\section{Introduction and Physics Motivation}
% ==============================================================================

The energy-momentum tensor (EMT) $T_{\mu\nu}(x)$ is a fundamental operator in quantum field theory. In QCD it encodes the distribution of energy, momentum, pressure, and shear stress inside hadrons. Its matrix elements between hadron states,

\begin{equation}
  \bra{\pi(p_f)} T_{\mu\nu} \ket{\pi(p_i)},
\end{equation}

determine the gravitational form factors $A(q^2)$, $B(q^2)$, $D(q^2)$, and $\bar{c}(q^2)$ of the pion. These form factors provide access to the mass and spin decomposition of the pion, its mechanical radius, and the $D$-term (pressure distribution). Unlike the proton, the pion has spin zero, simplifying the Lorentz decomposition.

On the lattice, $T_{\mu\nu}$ requires a non-perturbative renormalization scheme. The gradient flow \cite{Luscher:2010iy,Luscher:2013vga} provides a regularization-independent method: flowing gauge and fermion fields to finite flow time $t$ renders the bare composite operators finite. At small flow time, flowed operators are matched to $\overline{\text{MS}}$ via perturbatively computed matching coefficients $c_q(t)$ and $c_g(t)$:

\begin{equation}\label{eq:renorm_ansatz}
  T_{\mu\nu}^{\text{ren}}(x)
    = c_g(t)\, O_{\mu\nu}^g(t,x)
    + c_q(t)\, O_{\mu\nu}^q(t,x)
    + \text{mixing/trace terms}.
\end{equation}

The mixing must account for the fact that the trace of $T_{\mu\nu}$ receives contributions from the trace anomaly. The matching coefficients $c_q(t)$, $c_g(t)$ are computed in perturbation theory and are part of the analysis stage, not the contraction kernel.

This note documents the complete numerical implementation in PyQUDA of:
\begin{itemize}[nosep]
  \item Connected meson three-point functions with quark EMT insertion (Section~\ref{sec:connected}),
  \item Stochastic quark one-point functions for disconnected diagrams (Section~\ref{sec:quark1pt}),
  \item Gluon one-point functions (Section~\ref{sec:gluon1pt}),
\end{itemize}
together with the gradient-flow schedule and HDF5 output formats. The target audience is lattice practitioners who need a line-by-line correspondence between the continuum formulas and the einsum contraction patterns in the Python source.

\subsection*{Key conventions}
Throughout this note:
\begin{itemize}[nosep]
  \item Euclidean metric with $\gamma_4$ hermitian: $\gamma_\mu^\dagger = \gamma_\mu$, $\{\gamma_\mu,\gamma_\nu\} = 2\delta_{\mu\nu}$.
  \item $\gamma_5 = \gamma_1\gamma_2\gamma_3\gamma_4 = \gamma_5^\dagger$.
  \item Gamma hermiticity of the quark propagator: $S(x,y)^\dagger = \gamma_5 S(y,x) \gamma_5$.
  \item Convention B meson contraction with $\code{meson\_sign} = +1$.
  \item Source position $x_0 = (t_0, \vec{x}_0)$, insertion $x = (\tau, \vec{x})$, sink $y = (t_{\text{sep}}, \vec{y})$.
  \item Momentum phase: $\Phi_k(r - x_0) = e^{i k \cdot (r - x_0)}$ where $k$ uses continuum normalization $k_\mu \in \frac{2\pi}{L_\mu}\mathbb{Z}$.
\end{itemize}


% ==============================================================================
\section{Code Architecture Overview}
% ==============================================================================

The class hierarchy is:
\begin{center}
\begin{tabular}{ll}
  \code{EMTDisconnectedQuark1pt} & shared quark-1pt loop \\
  \code{EMTDisconnectedGluon1pt} & shared gluon-1pt loop \\
  \code{QuarkEMT(EMTDisconnectedQuark1pt)} & connected meson 3pt + 2pt \\
\end{tabular}
\end{center}

\code{QuarkEMT} inherits all stochastic 1pt methods (including covariant derivatives, gamma5 hermiticity, and flow utilities) while adding meson-specific two-point and three-point contractions.

The application-level entry points are:
\begin{itemize}[nosep]
  \item \code{application/EMT\_meson/perlmutter/Pyquda\_EMT\_quark\_3pt.py} -- connected 3pt
  \item \code{application/EMT\_disconnected\_1pt/perlmutter/} -- shared quark/gluon loops
\end{itemize}

% ==============================================================================
\section{Part I: Connected Meson EMT}
\label{sec:connected}
% ==============================================================================

\subsection{Meson Two-Point Function}
\label{sec:2pt}

The physical connected meson two-point function with source interpolator $\Gamma_{\text{src}}$ and sink interpolator $\Gamma_{\text{sink}}$ is

\begin{equation}\label{eq:C2_def}
  C_2^{\Gamma_{\text{sink}}}(p, t)
    = \sum_{\vec{x}} \Phi_{-p}(x - x_0)\,
      \Tr_{sc}\Big[
        S_{\text{anti}}(x, x_0)\; \Gamma_{\text{sink}}\;
        S_q(x, x_0)\; \Gamma_{\text{src}}
      \Big],
\end{equation}

where the antiquark line is built using gamma5 hermiticity:

\begin{equation}\label{eq:S_anti}
  S_{\text{anti}}(x, x_0) = \gamma_5\, S_q(x, x_0)^\dagger\, \gamma_5.
\end{equation}

In the code (\code{contract\_meson\_2pt}), the independent backward propagator \code{prop\_bw} is used to form the antiquark line:

\begin{equation}\label{eq:bw_prop}
  \text{bw\_prop}(x) = \gamma_5\; \texttt{prop\_bw}(x, x_0)^\dagger\; \gamma_5,
\end{equation}

and the contraction is computed as

\begin{equation}\label{eq:C2_code}
  C_2 = \sum_{\vec{x}} \Phi_{-p}(x - x_0)\;
        \Tr_{sc}\big[
          \text{bw\_prop}(x)\; \Gamma_{\text{sink}}\;
          \texttt{prop\_fw}(x, x_0)\; \Gamma_{\text{src}}
        \big].
\end{equation}

The implementation scans all 16 gamma structures for the sink while keeping $\Gamma_{\text{src}}$ fixed. The 16 gamma matrices (in \code{my\_gammas}) are:

{
\small
\begin{center}
\begin{tabular}{ll}
  \code{"5"} = $\gamma_5$, & \code{"T"} = $\gamma_4$, \\
  \code{"T5"} = $\gamma_4\gamma_5$, & \code{"X"} = $\gamma_1$, \\
  \code{"X5"} = $\gamma_1\gamma_5$, & \code{"Y"} = $\gamma_2$, \\
  \code{"Y5"} = $\gamma_2\gamma_5$, & \code{"Z"} = $\gamma_3$, \\
  \code{"Z5"} = $\gamma_3\gamma_5$, & \code{"I"} = $\mathbf{1}$, \\
  \code{"SXT"} = $\sigma_{14}$, & \code{"SXY"} = $\sigma_{12}$, \\
  \code{"SXZ"} = $\sigma_{13}$, & \code{"SYT"} = $\sigma_{24}$, \\
  \code{"SYZ"} = $\sigma_{23}$, & \code{"SZT"} = $\sigma_{34}$,
\end{tabular}
\end{center}
}

The raw PyQUDA tensor masks obey
$S_{\mu\nu}=\frac{1}{2}[\gamma_\mu,\gamma_\nu]$ without an extra $i$; multiply
them by $i$ if the Hermitian convention
$\sigma_{\mu\nu}=\frac{i}{2}[\gamma_\mu,\gamma_\nu]$ is desired.  Likewise the
raw bit-mask labels satisfy
$Y5_{\rm phys}=-Y5_{\rm raw}$ and
$T5_{\rm phys}=-T5_{\rm raw}$ when every physical axial matrix is defined as
$\gamma_\mu\gamma_5$.  These transformations are stored explicitly in
\code{physical\_from\_pyquda}. The ordering is fixed by \code{pyquda\_gammas\_order}:
\code{[15, 8, 7, 1, 14, 2, 13, 4, 11, 0, 9, 3, 5, 10, 6, 12]}.

In the einsum contraction, the data layout is
\begin{center}
\code{[parity, t, z, y, x, spin\_src, color\_src, spin\_sink, color\_sink]}
\end{center}
so the einsum pattern for forming $\text{bw\_prop}$ is
\begin{verbatim}
  bw_prop = contract("ij, wtzyxilab, kl -> wtzyxkjba",
                     G5, prop_bw.conj(), G5)
\end{verbatim}
where index \code{i,j,k,l} $\in \{0,1,2,3\}$ (spin) and \code{a,b} $\in \{0,1,2\}$ (color). The transpose of spin-sink and spin-source indices $(\texttt{ilab} \to \texttt{kjba})$ implements the dagger with spin-color reordering.

The final scalar folding into $C_2$ reads
\begin{verbatim}
  bw_prop = contract("wtzyxjicf, gim -> gwtzyxjmcf", bw_prop, sink_gammas)
  scalar   = contract("gwtzyxjiab, wtzyxilba, lj -> gwtzyx",
                      bw_prop, prop_fw, src_gamma)
\end{verbatim}
where the first line applies all 16 sink gamma matrices via a stacked index \code{g}, and the second line traces over spin-sink (\code{i}), spin-source (\code{l,j}), and color ($\Tr_c$ via repeated \code{a,b}).

\textbf{Boosted smearing:} If \code{CG\_GaussSmear} is enabled, Gaussian smearing with a momentum boost is applied at both source and sink:
\begin{itemize}[nosep]
  \item Source smearing uses boost \code{pos\_boost} on the forward source and \code{neg\_boost} on the backward source (when they differ).
  \item Sink smearing is applied during the 2pt contraction with the same \code{pos\_boost} and \code{neg\_boost}.
\end{itemize}
The smearing operator used throughout this note is the FFT boosted Gaussian
operator \(\mathcal{S}_{w,\mathbf{k}}\).  Its real-space kernel is
\begin{equation}\label{eq:emt_smear_kernel}
  K_{\mathbf{k}}(\mathbf{r})
  =
  \exp\!\left(-\frac{r_x^2+r_y^2+r_z^2}{2w^2}\right)
  \exp\!\left[
    2\pi i
    \left(
      \frac{k_x r_x}{L_x}
      + \frac{k_y r_y}{L_y}
      + \frac{k_z r_z}{L_z}
    \right)
  \right],
\end{equation}
where \(\mathbf{k}\) is the integer vector supplied by \code{pos\_boost} or
\code{neg\_boost}.  The current \code{boosted\_smearing\_pyquda.py}
implementation applies a three-dimensional spatial FFT convolution,
\begin{equation}\label{eq:emt_fft_smear}
  \psi^{\,{\rm smear}}_{\mathbf{k}}(\mathbf{x},t)
  =
  \mathcal{F}^{-1}_{3d}
  \left[
    \mathcal{F}_{3d}[\psi](\mathbf{p},t)\,
    \mathcal{F}_{3d}[K_{\mathbf{k}}](\mathbf{p})
  \right](\mathbf{x},t),
\end{equation}
to every spin-color source column of a propagator.  The FFT is spatial only and
the same kernel is used on each Euclidean time slice.  No explicit gauge
rotation is applied in the current implementation.

\subsection{Connected Three-Point Function (before sequential trick)}
\label{sec:3pt_pre_seq}

For a local quark bilinear insertion $\mathcal{O}_g(x)$ at the space-time point $x = (\tau, \vec{x})$ and a final sink with momentum $p_f$, the connected three-point function is

\begin{equation}\label{eq:C3_pre_seq}
  C_3^g(q, \tau; p_f, t_{\text{sep}})
    = \sum_{\vec{x}} \sum_{\vec{y}}
      \Phi_q(x - x_0)\; \Phi_{p_f}(y - x_0)\;
      \Tr_{sc}\Big[
        S_{\text{anti}}(y, x_0)\; \Gamma_{\text{sink}}\;
        S_q(y, x)\; \mathcal{O}_g(x)\; S_q(x, x_0)\; \Gamma_{\text{src}}
      \Big].
\end{equation}

Here $q$ is the momentum transfer (denoted \code{qext} in the code), $p_f$ is the fixed sink momentum (\code{pf}), and $t_{\text{sep}}$ is the source-sink separation. The notation follows the usual three-point pattern: the quark propagates from source $x_0$ to insertion $x$, the operator is inserted, the quark continues to sink $y$, and the antiquark runs directly from source to sink.

\subsection{Scalar Diagnostic: $C_3^\chi$}
\label{sec:C3_chi}

With $\mathcal{O}_g(x) = \mathbf{1}$ (the unit operator in spin-color space), the three-point function reduces to

\begin{equation}\label{eq:C3_chi_def}
  C_3^\chi(q, \tau)
    = \sum_{\vec{x}} \Phi_q(x - x_0)\;
      \Tr_{sc}\Big[
        S_{\text{seq\_anti}}(x)\;
        S_q(x, x_0)\;
        \Gamma_{\text{src}}
      \Big],
\end{equation}

where $S_{\text{seq\_anti}}(x)$ is the backward sequential meson line defined in Section~\ref{sec:seq_src}. This diagnostic verifies that the sequential source construction is correct: a well-constructed sequential source should give $C_3^\chi$ that is numerically consistent with the two-point function when the scalar insertion is trivial.

In \code{get\_C3\_chi}, the contraction is
\begin{verbatim}
  dst2 = make_dst2(seq_bw_prop)  # gamma5 * S_seq^dag * gamma5
  scalar_field = contract("wtzyxabij, wtzyxbcji, ca -> wtzyx",
                          dst2, prop_fw.data, src_gamma)
  # Fourier transform to momentum space, keep time
  slice_t = gatherLattice( contract("qwtzyx, wtzyx -> qt",
                                    phases_3pt, scalar_field) )
\end{verbatim}
Note the spin indices: \code{dst2} has layout \code{(..., a=spin\_sink, b=spin\_src, i=color\_sink, j=color\_src)}, while \code{prop\_fw.data} has \code{(..., b=spin\_src, c=spin\_sink, j=color\_src, i=color\_sink)}. The contraction \code{abij, bcji, ca} traces over spin and color, with \code{src\_gamma} providing the source interpolator.

\subsection{Connected Quark EMT Insertion}
\label{sec:quark_emt_connected}

The Euclidean connected quark EMT insertion is

\begin{equation}\label{eq:Tmunu_q_bilinear}
  T_{\mu\nu}^q(x) = \frac{1}{2}\Big[
    \gamma_\nu D_\mu - \overleftarrow{D}_\mu \gamma_\nu
  \Big],
\end{equation}

where $D_\mu$ is the symmetric covariant derivative acting to the right and $\overleftarrow{D}_\mu$ acts to the left. This bilinear is not Hermitian by itself; the symmetrization $\mu \leftrightarrow \nu$ is performed after contraction to obtain a Hermitian EMT.

When inserted into the three-point function, two distinct contributions arise depending on which quark line the derivative acts upon:

\subsubsection{First term: derivative on the forward (source-to-insertion) line}

\begin{equation}\label{eq:C3_first}
  C_3^{\mu\nu,\,q,\,\text{first}}(q, \tau)
    = +\frac{1}{2} \sum_{\vec{x}} \Phi_q(x - x_0)\;
      \Tr_{sc}\Big[
        S_{\text{seq\_anti}}(x)\;
        \gamma_\nu\; \big(D_\mu S_q(x, x_0)\big)\;
        \Gamma_{\text{src}}
      \Big].
\end{equation}

The derivative acts on the forward propagator \emph{before} multiplication by the gamma matrix. The overall sign is $+1/2$.

\subsubsection{Second term: derivative on the backward (sequential) line}

\begin{equation}\label{eq:C3_second}
  C_3^{\mu\nu,\,q,\,\text{second}}(q, \tau)
    = -\frac{1}{2} \sum_{\vec{x}} \Phi_q(x - x_0)\;
      \Tr_{sc}\Big[
        \big(\overleftarrow{D}_\mu S_{\text{seq\_anti}}(x)\big)\;
        \gamma_\nu\; S_q(x, x_0)\;
        \Gamma_{\text{src}}
      \Big].
\end{equation}

The derivative acts on the sequential backward propagator \emph{from the left}. The overall sign is $-1/2$, which follows from the minus sign in Eq.~\eqref{eq:Tmunu_q_bilinear} for the left-derivative term.

\subsubsection{Final contraction and symmetrization}

The full quark EMT three-point function is the sum of the two contributions:

\begin{equation}\label{eq:C3_Tmunu_sum}
  C_3^{\mu\nu,\,q}(q, \tau)
    = C_3^{\mu\nu,\,q,\,\text{first}}(q, \tau)
    + C_3^{\mu\nu,\,q,\,\text{second}}(q, \tau).
\end{equation}

After this sum, explicit symmetrization enforces $T_{\mu\nu} = T_{\nu\mu}$:

\begin{equation}\label{eq:symmetrization}
  C_3^{\mu\nu,\,q}(q, \tau) \to
    \frac{1}{2}\big(
      C_3^{\mu\nu,\,q}(q, \tau) + C_3^{\nu\mu,\,q}(q, \tau)
    \big).
\end{equation}

The code in \code{get\_C3\_Tmunu\_symmetrized} implements this as:

\begin{verbatim}
# First term: +1/2 * Tr[dst2 * gamma_nu * (D_mu prop_fw) * Gamma_src]
for mu in range(4):
    D_fw = covdev_sym_prop(U_f, prop_fw, mu)
    for nu in range(4):
        gamma_D_fw = contract("ab, wtzyxbdij -> wtzyxadij",
                              D_gammas[nu], D_fw.data)
        scalar = 0.5 * contract("wtzyxabij, wtzyxbcji, ca -> wtzyx",
                                dst2, gamma_D_fw, src_gamma)

# Second term: -1/2 * Tr[(leftD_mu dst2) * gamma_nu * prop_fw * Gamma_src]
for mu in range(4):
    leftD_dst2 = left_covdev_dst2_from_prop(U_f, seq_bw_prop, mu)
    for nu in range(4):
        gamma_fw = contract("ab, wtzyxbdij -> wtzyxadij",
                            D_gammas[nu], prop_fw.data)
        scalar = -0.5 * contract("wtzyxabij, wtzyxbcji, ca -> wtzyx",
                                 leftD_dst2, gamma_fw, src_gamma)

# Symmetrize
for mu in range(4):
    for nu in range(mu+1, 4):
        C3_Tmunu[:, mu, nu] = 0.5 * (C3_Tmunu[:, mu, nu]
                                    + C3_Tmunu[:, nu, mu])
        C3_Tmunu[:, nu, mu] = C3_Tmunu[:, mu, nu]
\end{verbatim}

The production contraction now evaluates all 16 gamma matrices for each of the
four derivative directions.  The primitive array is unsymmetrized; selecting
the vector channels $\gamma_1,\gamma_2,\gamma_3,\gamma_4$ and symmetrizing
reproduces the previous \code{C3\_Tmunu} exactly.  Here derivative index zero is
$x$, one is $y$, two is $z$, and three is Euclidean time.  The batched Gamma
contraction adds no propagator inversion, fermion-flow step, or covariant
derivative call.

\begin{remark}[Trace ordering]
The einsum pattern \code{abij, bcji, ca} should be read as:
\begin{itemize}[nosep]
  \item \code{a}: spin-sink of \code{dst2} (or \code{leftD\_dst2})
  \item \code{b}: spin-source of \code{dst2}, spin-sink of forward propagator
  \item \code{c}: spin-source of forward propagator, spin index of \code{src\_gamma}
  \item \code{i,j}: color-sink and color-source, traced by Einstein summation
  \item \code{ca}: \code{src\_gamma} with indices (spin\_source, spin\_source\_of\_gamma)
\end{itemize}
The color trace is implicit in the repeated \code{i} and \code{j} indices.
\end{remark}

\subsection{Fixed-Sink Sequential Source Construction}
\label{sec:seq_src}

The fixed-sink method absorbs the sum over the sink position $y$ into a sequential propagator. This avoids an $O(V^2)$ summation and makes the three-point calculation tractable.

\textbf{Step 1: Build the source-smeared, sink-smeared forward propagator.}

Starting from the source-smeared point-sink forward propagator $\mathtt{prop\_fw\_SP}(y, x_0)$, optional sink smearing is applied:

\begin{equation}
  \mathtt{prop\_fw\_SS}(y, x_0)
    = \text{boosted\_smearing}\big(
        \mathtt{prop\_fw\_SP}(\cdot, x_0),\;
        \text{boost} = \mathtt{pos\_boost}
      \big)(y).
\end{equation}

\textbf{Step 2: Restrict to the sink time slice and apply momentum phase.}

The function \code{create\_meson\_bw\_seq\_pyquda} (in \code{bw\_seq\_pyquda.py}) constructs the sequential source:

\begin{equation}\label{eq:eta_seq}
  \eta_{\text{seq}}(y; p_f, t_{\text{sep}})
    = \delta_{t_y,\, t_0 + t_{\text{sep}}}\;
      \Phi_{p_f}(y - x_0)\;
      \Gamma_{\text{seq}}\;
      \mathtt{prop\_fw\_SS}(y, x_0),
\end{equation}

where the time-slice restriction is implemented via \code{pyquda\_utils.source.sequential12}, and

\begin{equation}\label{eq:Gamma_seq}
  \Gamma_{\text{seq}} = \gamma_5\, \Gamma_{\text{sink}}^\dagger\, \gamma_5.
\end{equation}

In the einsum implementation:
\begin{verbatim}
  gamma_seq = einsum("ab, cb, cd -> ad", G5, sink_gamma.conj(), G5)
  seq_data = einsum("wtzyx, wtzyxilab -> wtzyxilab", mom_phase, seq_data)
  seq_data = einsum("ij, wtzyxjlab -> wtzyxilab", gamma_seq, seq_data)
\end{verbatim}

\textbf{Step 2b: Smear the active sequential source.}

For a smeared-smeared connected three-point function the code smears the
sequential source before inversion:
\begin{equation}\label{eq:eta_seq_smeared}
  \eta_{\text{seq}}^{SS}
  =
  \mathcal{S}_{w,\mathtt{neg\_boost}}
  \left[
    \delta_{t_y,\, t_0 + t_{\text{sep}}}\;
    \Phi_{p_f}(y - x_0)\;
    \Gamma_{\text{seq}}\;
    \mathtt{prop\_fw\_SS}(y, x_0)
  \right].
\end{equation}
The input \(\mathtt{prop\_fw\_SS}\) has already received the
\code{pos\_boost} sink smearing.  The outer
\(\mathcal{S}_{w,\mathtt{neg\_boost}}\) represents the active line associated
with the backward sequential propagator in the current pion EMT implementation.

\textbf{Step 3: Invert the Dirac operator.}

\begin{equation}
  D\; S_{\text{seq}} = \eta_{\text{seq}}.
\end{equation}

\textbf{Step 4: Form the backward sequential meson line.}

Using gamma5 hermiticity, the line that connects the insertion to the sink is

\begin{equation}\label{eq:S_seq_anti}
  S_{\text{seq\_anti}}(x; p_f, t_{\text{sep}})
    = \gamma_5\; S_{\text{seq}}(x)^\dagger\; \gamma_5.
\end{equation}

This is computed in \code{\_make\_dst2}:
\begin{verbatim}
  dst2 = contract("ab, wtzyxbcij, cd -> wtzyxadij",
                  G5, prop.data.conj().transpose(0,1,2,3,4,6,5,8,7), G5)
\end{verbatim}
The \code{transpose(0,1,2,3,4,6,5,8,7)} swaps spin-source/sink (indices 5$\leftrightarrow$6) and color-source/sink (indices 7$\leftrightarrow$8), implementing the Hermitian conjugate with the correct index layout.

\textbf{Step 5: Post-sequential three-point function.}

After the sequential inversion, the sink-summed three-point function takes the compact form:

\begin{equation}\label{eq:C3_post_seq}
  C_3^g(q, \tau; p_f, t_{\text{sep}})
    = \sum_{\vec{x}} \Phi_q(x - x_0)\;
      \Tr_{sc}\Big[
        S_{\text{seq\_anti}}(x; p_f, t_{\text{sep}})\;
        \mathcal{O}_g(x)\;
        S_q(x, x_0)\;
        \Gamma_{\text{src}}
      \Big].
\end{equation}

This is the master formula used by both \code{get\_C3\_chi} and \code{get\_C3\_Tmunu\_symmetrized}.


\subsection{Symmetric Covariant Derivative on a Propagator}
\label{sec:covdev_sym}

The symmetric (improved) covariant derivative acting on a quark propagator is defined on the lattice as

\begin{equation}\label{eq:D_sym}
  D_\mu^{\text{sym}} S(x)
    = \frac{1}{2}\Big[
        U_\mu(x)\, S(x + \hat{\mu})
        - U_{-\mu}(x)\, S(x - \hat{\mu})
      \Big],
\end{equation}

where $U_{-\mu}(x) = U_\mu^\dagger(x - \hat{\mu})$. The forward direction $+\mu$ corresponds to direction index $\mu \in \{0,1,2,3\}$, and the backward direction $-\mu$ corresponds to $\mu+4 \in \{4,5,6,7\}$ in the QUDA direction enumeration.

The code (\code{\_covdev\_sym\_prop}) applies this spin-color component by component via \code{MultiFermion}:

\begin{verbatim}
  mf = propagatorToMultiFermion(prop)
  for spin in range(4):
      for color in range(3):
          idx = spin * 3 + color
          Dp = U.pure_gauge.covDev(mf[idx], mu)      # +mu
          Dm = U.pure_gauge.covDev(mf[idx], mu + 4)   # -mu
          mf_covdev[idx] = 0.5 * (Dp - Dm)
\end{verbatim}

\begin{remark}[Sign convention]
The minus sign between $D_{+}$ and $D_{-}$ in the code, $0.5 \times (D_{+} - D_{-})$, corresponds to the standard lattice symmetric derivative where covDev implements the forward/backward parallel transport without the link division by $2a$. The factor $1/2$ in the formula provides the continuum normalization.
\end{remark}

\subsection{Left Derivative on the Sequential Line}
\label{sec:left_covdev}

The second term of the quark EMT insertion requires the left-acting derivative on the backward sequential line $S_{\text{seq\_anti}}$. Rather than implementing a separate left-derivative routine, the code exploits gamma5 hermiticity:

\begin{equation}\label{eq:leftD_via_gamma5}
  \overleftarrow{D}_\mu S_{\text{seq\_anti}}(x)
    = \gamma_5\,
      \big(D_\mu S_{\text{seq}}(x)\big)^\dagger\,
      \gamma_5.
\end{equation}

The derivation is:
\begin{align}
  \overleftarrow{D}_\mu S_{\text{seq\_anti}}
    &= \overleftarrow{D}_\mu \big( \gamma_5 S_{\text{seq}}^\dagger \gamma_5 \big) \\
    &= \gamma_5 \big( D_\mu S_{\text{seq}} \big)^\dagger \gamma_5,
\end{align}
where the derivative passes through the constant gamma matrices and acts on $S_{\text{seq}}^\dagger$ from the right, which is equivalent to the Hermitian conjugate of the right-acting derivative.

The implementation in \code{\_left\_covdev\_dst2\_from\_prop} is:

\begin{verbatim}
  D_y = covdev_sym_prop(U_f, prop, mu)   # D_mu S_seq
  D_y_dag = D_y.data.conj().transpose(0,1,2,3,4,6,5,8,7)  # (D_mu S_seq)^dag
  leftD_dst2 = contract("ab, wtzyxbcij, cd -> wtzyxadij",
                        G5, D_y_dag, G5)
\end{verbatim}

The resulting \code{leftD\_dst2} encodes $\overleftarrow{D}_\mu S_{\text{seq\_anti}}(x)$ and is used directly in the second-term contraction of Eq.~\eqref{eq:C3_second}.

\subsection{Gradient Flow}
\label{sec:flow_connected}

Gradient flow is applied to smooth UV fluctuations. The flow radius is approximately $\sqrt{8t}$ where $t$ is the flow time. The code supports both Wilson flow and Symanzik (L\"uscher-Weisz) flow via the \code{flow\_type} parameter.

\subsubsection{Flow schedule}

The flow schedule follows a \textbf{measure-before-flow} convention:

\begin{enumerate}[nosep]
  \item \textbf{Step 0}: Measure observables on the unflowed fields. Output index \code{step=0} is unflowed.
  \item \textbf{First interval (step 0 $\to$ 1)}: Advance fields by 10 small sub-steps of size $\varepsilon/10$.
  \item \textbf{Subsequent intervals}: Advance by one step of size $\varepsilon$.
\end{enumerate}

Thus output index \code{step} corresponds to flow time $t = \text{step} \times \varepsilon$.

\subsubsection{Flowing propagators together}

The forward propagator and sequential backward propagator are flowed simultaneously as a single \code{MultiField} on the flowed gauge background:

\begin{verbatim}
  mf_a = propagatorToMultiFermion(prop_a)
  mf_b = propagatorToMultiFermion(prop_b)
  packed = multiField(mf_a_all + mf_b_all)
  packed_flow = U_f.gradientFlow(packed, flow_type, Nsteps, stepsize)
\end{verbatim}

This ensures both fermion fields see the same flowed gauge links at every step.

\subsubsection{Gauge field flow}

The gauge field itself is flowed by \code{U\_f.wilsonFlow()} or \code{U\_f.symanzikFlow()}, which modifies the gauge links in place. Before flowing, anti-periodic boundary conditions in time are set via \code{U\_f.setAntiPeriodicT()}.

\subsection{Connected 3pt High-Level Algorithm}
\label{sec:connected_algorithm}

The \code{connected\_3pt} method orchestrates the full measurement:

\begin{enumerate}[nosep]
  \item Build source-smeared point-sink forward and backward propagators (\code{\_make\_meson\_source\_props}).
  \item Contract the meson 2pt function as a diagnostic (\code{contract\_meson\_2pt}).
  \item For each source-sink separation $t_{\text{sep}}$:
  \begin{enumerate}
    \item Build and invert the meson fixed-sink sequential source (\code{create\_meson\_bw\_seq\_pyquda}).
    \item For each flow step:
    \begin{enumerate}
      \item Compute $C_3^\chi(q,\tau)$ via \code{get\_C3\_chi}.
      \item Compute $C_3^{\mu\nu,q}(q,\tau)$ via \code{get\_C3\_Tmunu\_symmetrized}.
      \item Advance the flowed propagators via \code{\_advance\_flowed\_props}.
    \end{enumerate}
  \end{enumerate}
  \item Save the results to HDF5.
\end{enumerate}

\subsection{HDF5 Output Format (Connected 3pt)}
\label{sec:hdf5_3pt}

The output file (written by \code{save\_emt\_quark\_3pt\_hdf5}) contains:

\begin{itemize}[nosep]
  \item \code{C3\_chi}: 4D array \code{[N\_ts, Nsteps+1, Nq, Nt]} -- scalar diagnostic.
  \item \code{C3\_Tmunu}: 6D array \code{[N\_ts, Nsteps+1, Nq, 4, 4, Nt]} -- quark EMT 3pt.
  \item \code{C3\_local\_bilinear}: \code{[N\_ts,16,Nq,Nsteps+1,Nt]}.
  \item \code{C3\_derivative\_bilinear}: \code{[N\_ts,16,4,Nq,Nsteps+1,Nt]}.
  \item \code{gamma\_list}, \code{gamma\_pyquda\_ids},
        \code{gamma\_matrices}, \code{physical\_from\_pyquda}, and
        \code{derivative\_directions}: exact operator-basis metadata.
  \item \code{momentum\_transfer\_list}: list of $[q_x, q_y, q_z, q_t]$ arrays.
  \item HDF5 attributes: \code{measurement}, \code{spin}, \code{flow\_type}, \code{flow\_epsilon}, \code{flow\_steps}, \code{n\_t\_separations}, \code{src\_interpolator}, \code{sink\_interpolator}, \code{contraction\_convention} (``B''), \code{meson\_sign} ($+1$), \code{n\_qext}.
\end{itemize}

The meson 2pt diagnostic is saved separately by \code{save\_emt\_meson\_2pt\_hdf5} with:
\begin{itemize}[nosep]
  \item \code{C2}: 3D array \code{[16, N\_mom, Nt]} -- 16 sink gammas $\times$ all momenta $\times$ time.
  \item \code{gamma\_list}: list of gamma labels.
  \item \code{momentum\_list}: list of momentum vectors.
\end{itemize}


% ==============================================================================
\section{Part II: Disconnected One-Point Functions}
% ==============================================================================

\subsection{Stochastic Quark One-Point Function}
\label{sec:quark1pt}

The disconnected quark loop requires the trace of the quark propagator with operator insertions. These are estimated stochastically using $\znum$ noise vectors.

\subsubsection{Noise generation}

A $\znum$ noise fermion $\xi(x)$ is generated with random complex phases:

\begin{equation}\label{eq:zn_noise}
  \xi(x) \in \big\{ e^{2\pi i r / n} \mid r \in \{0, 1, \dots, n-1\} \big\},
\end{equation}

with \code{n = n\_zn}; the production default is $n=4$.  Each phase is
constructed by the fixed counter described in
\code{docs/EMT\_disconnected\_1pt/EMT\_disconnected\_1pt.tex}, using the
global coordinate, spin, color, configuration, base-noise index, and stream
salt.  Thus the global source is unchanged by the MPI decomposition and obeys
$\langle \xi(x) \rangle = 0$ and
$\langle \xi(x) \xi^\dagger(y) \rangle = \delta_{xy}$ in the noise ensemble.

The solution vector is obtained by inverting the Dirac operator:

\begin{equation}\label{eq:eta_noise}
  \eta(x) = D^{-1} \xi(x).
\end{equation}

\subsubsection{Identity channel and flowed-noise norm}

Two scalar quantities monitor the stochastic estimator:

\begin{align}
  \text{CHI}[0](q, t) &= \sum_{\vec{x}} \Phi_q(x)\;
    \xi^\dagger(x)\, \eta(x)
    \label{eq:CHI0}
    \\
  \text{CHI}[1](q, t) &= \sum_{\vec{x}} \Phi_q(x)\;
    \xi^\dagger(x)\, \xi(x)
    \label{eq:CHI1}
\end{align}

The first quantity is the identity channel of the complete 16-Gamma local
basis and is not duplicated.  The second is the flowed-noise norm.

In code:
\begin{verbatim}
  scalar_trace = local_bilinear[gamma_list.index("I")]

  dot_xi_xi = contract("etzyxbc, etzyxbc -> etzyx",
                        xi.data.conj(), xi.data)
  flowed_noise_norm = impose_P_Breit_slice(dot_xi_xi, phases_3pt)
\end{verbatim}

\subsubsection{Quark EMT one-point building block}

The stochastic estimate of the quark EMT one-point function is

\begin{equation}\label{eq:Tmunu_1pt_stoch}
  T_{\nu\mu}^q(q, t)
    = -\frac{1}{2} \sum_{\vec{x}} \Phi_q(x)\;
      \xi^\dagger(x)\;
      \gamma_\nu\,
      \big[ D_{+\mu} - D_{-\mu} \big]\,
      \eta(x).
\end{equation}

Here $D_{\pm\mu}$ are the forward/backward covariant derivatives (NOT the symmetrized versions -- the difference $D_+ - D_-$ already implements the symmetric derivative up to a factor of 2). The indices on $T_{\nu\mu}^q$ are ordered $(\nu, \mu)$ because the gamma matrix index comes first in the code:

\begin{verbatim}
  for mu in range(4):
      tmp = covDev(eta, mu) - covDev(eta, mu + 4)   # (D_+mu - D_-mu) eta
      for nu in range(4):
          Y = contract("ab, ...bc -> ...ac", D_gammas[nu], tmp.data)
          complex_field = contract("...sc, ...sc -> ...",
                                   xi.data.conj(), Y)
          Tmunu[nu, mu] += -0.5 * impose_P_Breit_slice(...)
\end{verbatim}

The minus sign ($-1/2$) follows from Eq.~\eqref{eq:Tmunu_q_bilinear} combined with the stochastic trace convention.

\subsubsection{Symmetrization, noise averaging, and volume normalization}

After accumulating all $\mu,\nu$ components, the tensor is symmetrized (identical to the connected case):

\begin{equation}
  T_{\mu\nu}^q \to \frac{1}{2}\big( T_{\mu\nu}^q + T_{\nu\mu}^q \big).
\end{equation}

The measurements are repeated for \code{n\_vec} independent noise vectors and averaged:

\begin{equation}
  \langle T_{\mu\nu}^q \rangle
    = \frac{1}{N_{\text{vec}}} \sum_{v=1}^{N_{\text{vec}}}
      T_{\mu\nu}^{q,(v)}.
\end{equation}

Finally, the result is divided by the spatial volume $V_3 = L_x L_y L_z$ to obtain the intensive one-point density:

\begin{equation}
  T_{\mu\nu}^{q,\,\text{avg}}(q, t, t_f)
    = \frac{1}{V_3} \langle T_{\mu\nu}^q (q, t, t_f) \rangle.
\end{equation}

\subsubsection{HDF5 output format (Quark 1pt)}

The canonical EMTc (written by the shard finalizer) contains:

\begin{itemize}[nosep]
  \item \code{raw/local\_bilinear\_pervec}: \code{[n\_vec,16,Nq,Nsteps+1,Nt]}.
  \item \code{raw/derivative\_bilinear\_pervec}: \code{[n\_vec,16,4,Nq,Nsteps+1,Nt]}.
  \item \code{raw/flowed\_noise\_norm\_pervec}: \code{[n\_vec,Nq,Nsteps+1,Nt]}.
  \item \code{avg/local\_bilinear} and \code{avg/derivative\_bilinear}: source-averaged, volume-normalized primitives.
  \item \code{avg/Tmunu/T11, T12, ..., T44}: averaged, volume-normalized tensor components (symmetric storage: only $\mu \leq \nu$).
  \item \code{avg/flowed\_noise\_norm}: \code{[Nq,Nsteps+1,Nt]}.
  \item Attributes: \code{measurement}, \code{flow\_type}, \code{flow\_epsilon}, \code{flow\_steps}, \code{n\_vec}, \code{n\_zn}.
\end{itemize}

\subsection{Ringed Fermion Normalization}
\label{sec:ringed_fermion}

The ringed fermion renormalization scheme \cite{Makino:2016} requires the
flowed fermion kinetic expectation value.  At zero momentum, the raw diagonal
quark EMT one-point gives, for each effective stochastic source \(r\),

\begin{align}\label{eq:ringed_norm}
  K_r(t_f,\tau)
    &= -\frac{2}{V_s}\sum_{\mu=1}^4
       T_{\mu\mu,r}^q(0,t_f,\tau) \\
    &= \frac{1}{V_s}\sum_{\mu,\vec{x}}
       \xi^\dagger(x,t_f)\;
       \gamma_\mu\,
       \big[ D_{+\mu} - D_{-\mu} \big]\,
       \eta(x,t_f).
\end{align}

The EMT 1pt run extracts this identity directly from
the vector diagonal channels of
\code{raw/derivative\_bilinear\_pervec}, without another solve, flow update, derivative, or
gather, and writes \code{derived/ringed/kinetic\_pervec} and
\code{derived/ringed/kinetic\_spacetime} in the same EMTc; it does not contain
configuration-local \code{Z\_ring} datasets.  The same quantity can be checked
against

\begin{center}
\code{avg/Tmunu/T11}, \code{T22}, \code{T33}, \code{T44}
\end{center}

at the \(q=0\) momentum index.  Final ringed factors must be evaluated from the
ensemble average of \code{derived/ringed/kinetic\_spacetime}; averaging factors formed as
the inverse of each configuration's kinetic value is not equivalent.  The
local identity and flowed-noise norm remain diagnostics.  The complete
convention and schema are documented in the disconnected one-point note.
Production one-point jobs write base/HP-interval shards by default.  Only the
explicit finalizer publishes the canonical EMTc with embedded kinetic data after
all configured bases have complete HP coverage.

\subsection{Gluon One-Point Function}
\label{sec:gluon1pt}

\subsubsection{Clover field strength}

The gauge field strength $F_{\mu\nu}(x)$ is constructed via the clover (plaquette + rectangle) operator. The standard clover loop for the $(\mu,\nu)$ plane combines four oriented paths:

\begin{equation}\label{eq:clover_loops}
  \begin{aligned}
    \text{loop}_1 &: \mu \to \nu \to -\mu \to -\nu, \\
    \text{loop}_2 &: \nu \to -\mu \to -\nu \to \mu, \\
    \text{loop}_3 &: -\mu \to -\nu \to \mu \to \nu, \\
    \text{loop}_4 &: -\nu \to \mu \to \nu \to -\mu.
  \end{aligned}
\end{equation}

The field strength is extracted as

\begin{equation}\label{eq:F_clover}
  F_{\mu\nu}(x) = \frac{1}{8}\Big(
    Q_{\mu\nu}(x) - Q_{\mu\nu}^\dagger(x)
  \Big),
\end{equation}

where $Q_{\mu\nu}$ is the sum of the four clover loops (each with unit coefficient). The resulting $F_{\mu\nu}$ is anti-Hermitian by construction.

\subsubsection{Traceless projection}

The field strength is projected onto the traceless anti-Hermitian $\mathfrak{su}(3)$ algebra:

\begin{equation}\label{eq:F_traceless}
  F_{\mu\nu}^{\text{su(3)}}(x)
    = F_{\mu\nu}(x) - \frac{1}{N_c} \Tr_c\big[F_{\mu\nu}(x)\big] \cdot \mathbf{1}_{N_c}.
\end{equation}

The code then multiplies by $-i$ to match the Euclidean convention:

\begin{equation}\label{eq:F_minus_i}
  F_{\mu\nu}^{\text{Eucl}}(x) = -i\, F_{\mu\nu}^{\text{su(3)}}(x).
\end{equation}

In \code{\_F\_clover\_traceless}:
\begin{verbatim}
  A = 0.125 * (data - data.swapaxes(-2, -1).conjugate())
  trA = contract("...ii -> ...", A)
  A -= trA[..., None, None] * I / Nc
  data[...] = (-1j) * A
\end{verbatim}

\subsubsection{Gluon EMT building block}

The gluonic EMT one-point function is

\begin{equation}\label{eq:Tmunu_gluon}
  T_{\mu\nu}^g(q, t)
    = \frac{2}{V_3} \sum_{\vec{x}} \Phi_q(x)
      \sum_{\substack{\rho=0\\ \rho \neq \mu,\nu}}^{3}
      \Tr_c\Big[
        F_{\mu\rho}(x)\; F_{\nu\rho}(x)
      \Big].
\end{equation}

The sum over $\rho$ runs over all Lorentz indices except $\mu$ and $\nu$. This is the full gluonic EMT building block (not a traceless projection of the EMT itself -- the traceless projection in \code{\_F\_clover\_traceless} operates only on the field-strength level to ensure $\mathfrak{su}(3)$ algebra).

In code:
\begin{verbatim}
  for mu in range(4):
      for nu in range(mu, 4):
          tmp = zeros(...)
          for rho in range(4):
              if rho == mu or rho == nu: continue
              tmp += contract("...ab, ...ba -> ...", F[mu][rho], F[nu][rho])
          Tmunu_t[mu, nu, :, step, :] += 2.0 * slice_t
\end{verbatim}

The factor 2 is explicit in the formula, and the volume normalization $1/V_3$ occurs after the spatial Fourier sum.

\textbf{Antisymmetry:} $F_{\mu\nu}$ is antisymmetric by construction, so only the six independent planes $(0,1), (0,2), (0,3), (1,2), (1,3), (2,3)$ are computed, and $F_{\nu\mu} = -F_{\mu\nu}$ is filled by antisymmetry.

\subsubsection{HDF5 output format (Gluon 1pt)}

The file (written by \code{save\_emt\_gluon\_1pt\_hdf5}) contains:

\begin{itemize}[nosep]
  \item \code{Tmunu/T11, T12, ..., T44}: 3D arrays \code{[Nq, Nsteps+1, Nt]} -- volume-normalized tensor components (symmetric storage: only $\mu \leq \nu$).
  \item Attributes: \code{measurement}, \code{flow\_type}, \code{flow\_epsilon}, \code{flow\_steps}.
\end{itemize}

\subsection{Gradient Flow Schedule (One-Point)}
\label{sec:flow_1pt}

The one-point flow schedule mirrors the connected 3pt schedule (measure-before-flow, step-0 is unflowed). There is one important difference for the quark 1pt: the noise vector $\xi$ and solution $\eta$ are flowed together as a single \code{MultiField}. This ensures that the flowed $\xi$ and $\eta$ correspond to the same flowed gauge background.

\textbf{Quark 1pt flow:}
\begin{verbatim}
  for step in range(Nsteps + 1):
      measure primitive bilinears, flowed_noise_norm
      if Nsteps > 0 and step == 0:
          temp = multiField([xi, eta])
          temp_flow = U_f.gradientFlow(temp, flow_type, 10, stepsize/10)
          xi, eta = temp_flow[0], temp_flow[1]
      elif Nsteps > 0 and step < Nsteps:
          temp = multiField([xi, eta])
          temp_flow = U_f.gradientFlow(temp, flow_type, 1, stepsize)
          xi, eta = temp_flow[0], temp_flow[1]
\end{verbatim}

\textbf{Gluon 1pt flow:}
\begin{verbatim}
  for step in range(Nsteps + 1):
      measure Tmunu
      if Nsteps > 0 and step == 0:
          U_flow.wilsonFlow(10, epsilon=stepsize/10)  # or symanzik
      elif Nsteps > 0 and step < Nsteps:
          U_flow.wilsonFlow(1, epsilon=stepsize)       # or symanzik
\end{verbatim}

\subsubsection{Full-volume source for a time-resolved flowed loop}

The quark one-point estimator retains the absolute insertion time even though
its initial stochastic source is four dimensional.  Let \(K(t_f)\) be the
gauge-covariant fermion-flow kernel and \(P_\tau\) the projector onto the
absolute insertion time \(\tau\).  The flowed stochastic fields and estimator
are
\begin{equation}
  \xi_f=K(t_f)\xi,
  \qquad
  \eta_f=K(t_f)D^{-1}\xi,
  \qquad
  \widehat L(\tau,t_f)
  =\xi^\dagger K^\dagger P_\tau\Gamma K D^{-1}\xi.
\end{equation}
Consequently,
\begin{equation}
  \mathbb{E}_\xi[\widehat L(\tau,t_f)]
  =\Tr\!\left[P_\tau\Gamma K(t_f)D^{-1}K^\dagger(t_f)\right].
  \label{eq:pion_full_volume_flowed_loop}
\end{equation}
The projector keeps a spatial trace at fixed \(\tau\); it does not sum over
insertion time.  Nevertheless, fermion flow uses the four-dimensional
covariant Laplacian and spreads in time with characteristic radius
approximately \(\sqrt{8t_f}\).  Restricting the initial source to one time
projector generally omits finite-flow contributions.  A complete time-dilution
basis would remain unbiased only after summing every projector, with the
corresponding inversion cost.  The current workflow instead uses full-volume
counter-based \(Z_4\) noise and stores all \(\tau\).  The complete derivation is
given in
\code{docs/EMT\_disconnected\_1pt/EMT\_disconnected\_1pt.tex}.

Spatial and four-dimensional hierarchical probing do not change the source
support.  Both multiply a full-volume base noise; spatial HP only makes the
probing sign independent of time.

\subsection{Disconnected Matrix Element Analysis}
\label{sec:analysis}

The contraction kernel outputs \emph{bare flowed operators} to HDF5 files (Sec.~\ref{sec:hdf5_summary}).
This section documents the full analysis chain that turns those outputs into physical gravitational form factors.
The procedure is organized in six steps.

% ------------------------------------------------------------------------------
\subsubsection{Step 1 -- Read input data per configuration}
\label{sec:analysis_step1}

For each gauge configuration, four HDF5 files are read:

\begin{center}
\begin{tabular}{@{}lll@{}}
\toprule
\textbf{File} & \textbf{Dataset} & \textbf{Shape} \\
\midrule
\code{EMT2pt/...h5}   & \code{C2}          & \code{[16, Np, Nt]} \\
                       & \code{gamma\_list} & 16 strings \\
                       & \code{momentum\_list} & \code{[Np, 4]} \\
\midrule
\code{EMT3pt/...h5}   & \code{C2}          & \code{[Nt]} (zero-momentum, one gamma) \\
                       & \code{C3\_chi}     & \code{[Nts, Ns+1, Nq, Nt]} \\
                       & \code{C3\_Tmunu}   & \code{[Nts, Ns+1, Nq, 4, 4, Nt]} \\
\midrule
\code{EMTc/...h5}     & \code{avg/Tmunu/Tij} & \code{[Nq, Ns+1, Nt]} each \\
                       & \code{avg/flowed\_noise\_norm} & \code{[Nq, Ns+1, Nt]} \\
                       & \code{raw/local\_bilinear\_pervec} & \code{[Nv,16,Nq,Ns+1,Nt]} \\
                       & \code{raw/derivative\_bilinear\_pervec} & \code{[Nv,16,4,Nq,Ns+1,Nt]} \\
\midrule
\code{EMTg/<lat>.EMTg.<cfg>.<ama>.<sm>.h5} & \code{Tmunu/Tij} & \code{[Nq, Ns+1, Nt]} each \\
\bottomrule
\end{tabular}
\end{center}

Notation: \code{Np} = number of 2pt momenta, \code{Nq} = number of momentum transfers,
\code{Nts} = number of $t_{\text{sep}}$ values, \code{Ns} = number of flow steps,
\code{Nv} = number of noise vectors, \code{Nt} = temporal lattice extent.

For the disconnected analysis, the relevant interpolator is $\Gamma = \gamma_5$ (pseudoscalar, labeled \code{"5"} in \code{gamma\_list}).
The sink-gamma index \code{ig5} is located by matching against \code{gamma\_list}.
The momentum index \code{ip} for a desired sink momentum $p_f$ is located by matching against \code{momentum\_list}.

% ------------------------------------------------------------------------------
\subsubsection{Step 2 -- Time alignment: source position and absolute vs.\ relative time}
\label{sec:analysis_step2}

The meson 2pt function $C_2$ is computed with a specific source located at
$x_0 = (t_0, \vec{x}_0)$.  The source time $t_0$ is encoded in the file tag
(e.g., \code{...x0y0z0t32...}).  The 2pt dataset \code{C2} has shape
\code{[16, Np, Nt]} where the last axis is the \textbf{absolute sink time}
$t_{\text{sink}} \in [0, N_t - 1]$.

The 1pt loops $L_{\mu\nu}(q, \tau)$ also use \textbf{absolute lattice time}
$\tau \in [0, N_t - 1]$ for the operator insertion.  The 1pt measurement is
source independent: its canonical EMTc file omits the \code{src} label and one
full-time loop file per configuration is reused for every two-point source
time.  The value of $t_0$ must therefore be read from the corresponding hadron
two-point file, not from the EMTc file.

The physically meaningful time coordinates are \emph{relative to the source}:

\begin{equation}\label{eq:time_alignment}
  t_{\text{sep}} = (t_{\text{sink}} - t_0) \bmod N_t, \qquad
  \tau_{\text{rel}} = (\tau_{\text{abs}} - t_0) \bmod N_t.
\end{equation}

In practice the source is placed away from the temporal boundary and
$t_{\text{sep}} \ll N_t$, so the modulo operation is a formality.
The analysis loops over $t_{\text{sep}}$ values of interest and, for each,
extracts $C_2$ at $t_{\text{sink}} = t_0 + t_{\text{sep}}$ and $L$ at
$\tau_{\text{abs}} = t_0 + \tau_{\text{rel}}$ (with trivial wrap-around handling).

\textbf{Location of $t_0$ in the HDF5 files.}
For source-dependent hadron correlators, the source position $x_0$ is stored in:
\begin{itemize}[nosep]
  \item The file name (e.g., \code{...x0y0z0t32...} -- the $t$ component is $t_0$).
  \item The HDF5 attribute \code{src\_interpolator} or \code{source\_position}
    (convention varies by file type; the hadron-correlator file name is the
    authoritative source).
\end{itemize}
The canonical quark loop name is
\code{EMTc/<lat>.EMTc.<cfg>.<ama>.<sm>.h5} and intentionally contains no
source position.

Parsing example (Python pseudocode):
\begin{verbatim}
  import re
  fname = "lat.EMT2pt.cfg100.ama0.x0y0z0t32.sm0.h5"
  t0 = int(re.search(r"t(-?\d+)", fname).group(1))
\end{verbatim}

% ------------------------------------------------------------------------------
\subsubsection{Step 3 -- Kinematics and momentum matching}
\label{sec:analysis_step3}

The meson 2pt and the 1pt loops have independent momentum lists.
The source pion carries momentum $p_i$, the sink carries $p_f$ (from
\code{momentum\_list} in the 2pt file), and the operator insertion carries
momentum $q$ (from \code{qext}).  Momentum conservation at the operator vertex
requires

\begin{equation}\label{eq:kinematics}
  q = p_f - p_i.
\end{equation}

The source momentum $p_i$ is determined by the source construction:
\begin{itemize}[nosep]
  \item For a point source without boost: $p_i = (0, \vec{0})$.
  \item For a Gaussian-smeared source with boost \code{pos\_boost}:
    $p_i = \mathtt{pos\_boost}$ (up to smearing-induced smearing of the
    momentum distribution).
\end{itemize}

Thus, for each pair $(p_f \in \mathtt{p\_2pt},\; q \in \mathtt{qext})$
satisfying $q = p_f - p_i$, the relevant 2pt and 1pt entries are matched.
The \textbf{simplest and most common case} is $p_i = 0$, for which
$p_f = q$ and the momentum lists can be chosen identical
($\mathtt{p\_2pt} = \mathtt{qext}$).

In the analysis loop, for each momentum transfer $q$:
\begin{enumerate}[nosep]
  \item Compute $p_f = q + p_i$.
  \item Find the index \code{ip} of $p_f$ in \code{momentum\_list}.
  \item Find the index \code{iq} of $q$ in \code{qext} (from HDF5 attributes
    or the \code{momentum\_transfer\_list} dataset).
  \item The 2pt at this kinematics is \code{C2[ig5, ip, :]}.
  \item The 1pt loop at this kinematics is \code{Tmunu[iq, :, :]}.
\end{enumerate}

% ------------------------------------------------------------------------------
\subsubsection{Step 4 -- Form the disconnected three-point function}
\label{sec:analysis_step4}

For a single gauge configuration, fixed flow step \code{is} (flow time
$t_f = \code{is} \times \varepsilon$), fixed Lorentz indices $(\mu,\nu)$,
and a matched momentum pair $(p_f, q)$, the disconnected three-point function
at relative insertion time $\tau$ and source-sink separation $t_{\text{sep}}$ is

\begin{equation}\label{eq:C3_disc_per_cfg}
  C_3^{\text{disc},\,q}(q, \tau, t_{\text{sep}}, t_f)\big|_{\text{cfg}}
    = C_2\big(p_f,\; t_{\text{sink}}{=}t_0{+}t_{\text{sep}}\big)\;
      \times\;
      L_{\mu\nu}^q\big(q,\; \tau_{\text{abs}}{=}t_0{+}\tau,\; t_f\big).
\end{equation}

Indices on the right-hand side refer to the stored array axes:
\begin{itemize}[nosep]
  \item \code{C2[ig5, ip, t0 + tsep]} -- 2pt at sink gamma $\gamma_5$,
    sink momentum $p_f$, absolute sink time $t_0 + t_{\text{sep}}$.
  \item \code{L\_quark[iq, istep, t0 + tau]} -- quark loop at momentum
    transfer $q$, flow step \code{is}, absolute insertion time $t_0 + \tau$.
  \item \code{L\_gluon[iq, istep, t0 + tau]} -- gluon loop, same indexing.
\end{itemize}

The gluon loop enters identically to the quark loop:

\begin{equation}\label{eq:C3_disc_gluon_per_cfg}
  C_3^{\text{disc},\,g}(q, \tau, t_{\text{sep}}, t_f)\big|_{\text{cfg}}
    = C_2\big(p_f,\; t_0{+}t_{\text{sep}}\big)\;
      \times\;
      L_{\mu\nu}^g\big(q,\; t_0{+}\tau,\; t_f\big).
\end{equation}

Only the window $0 < \tau < t_{\text{sep}}$ contributes to the ground-state
signal (the operator must be inserted between source and sink).
Points outside this window serve as cross-checks (should be consistent with
zero after vacuum subtraction).

The connected three-point function $C_3^{\text{conn}}$ is read directly from
\code{EMT3pt/...h5} (dataset \code{C3\_Tmunu}, shape
\code{[Nts, Ns+1, Nq, 4, 4, Nt]}).  It already includes the sequential-source
summation and uses the same source $x_0$; its time axis is absolute and must
be aligned identically.

% ------------------------------------------------------------------------------
\subsubsection{Step 5 -- Ensemble averaging and vacuum subtraction}
\label{sec:analysis_step5}

Let $N_{\text{cfg}}$ be the number of gauge configurations.  For each jackknife or bootstrap
resampling bin $b$, compute the ensemble averages:

\begin{align}
  \langle C_2(p, t_{\text{sep}}) \rangle_b
    &= \frac{1}{N_{\text{cfg}} - 1} \sum_{\text{cfg} \neq b}
       C_2(p, t_{\text{sep}})\big|_{\text{cfg}}, \label{eq:avg_C2} \\[4pt]
  \langle L_{\mu\nu}(q, \tau, t_f) \rangle_b
    &= \frac{1}{N_{\text{cfg}} - 1} \sum_{\text{cfg} \neq b}
       L_{\mu\nu}(q, \tau, t_f)\big|_{\text{cfg}}, \label{eq:avg_L} \\[4pt]
  \langle C_2 \cdot L_{\mu\nu} \rangle_b
    &= \frac{1}{N_{\text{cfg}} - 1} \sum_{\text{cfg} \neq b}
       \big[ C_2(p, t_{\text{sep}}) \; L_{\mu\nu}(q, \tau, t_f) \big]_{\text{cfg}}. \label{eq:avg_C2L}
\end{align}

The disconnected three-point function with vacuum subtraction is then

\begin{equation}\label{eq:C3_disc_ens}
  \boxed{
    C_3^{\text{disc}}(q, \tau, t_{\text{sep}}, t_f)
      = \langle C_2 \cdot L \rangle
      - \langle C_2 \rangle \; \langle L \rangle
  }.
\end{equation}

This subtraction removes the vacuum expectation value of the operator and the
disconnected contribution where the pion propagates without interacting with the EMT.
Eq.~\eqref{eq:C3_disc_ens} is applied separately to the quark loop ($L^q$) and gluon loop ($L^g$).

The jackknife error on $C_3^{\text{disc}}$ is computed in the standard way:
$\sigma^2 = \frac{N_{\text{cfg}} - 1}{N_{\text{cfg}}} \sum_b (C_3^{\text{disc}}(b) - \overline{C_3^{\text{disc}}})^2$.

% ------------------------------------------------------------------------------
\subsubsection{Step 6 -- Ratio method and ground-state extraction}
\label{sec:analysis_step6}

The bare matrix element is isolated by forming the ratio with the two-point function:

\begin{equation}\label{eq:ratio}
  R_{\mu\nu}^{\mathcal{O}}(q, \tau, t_{\text{sep}}, t_f)
    = \frac{C_3^{\mathcal{O}}(q, \tau, t_{\text{sep}}, t_f)}
           {\langle C_2(p_f,\, t_{\text{sep}}) \rangle},
\end{equation}

where $\mathcal{O} \in \{\text{conn},\, \text{disc-}q,\, \text{disc-}g\}$.
For the connected piece, $C_3^{\text{conn}}$ from the 3pt file enters the numerator.
For the disconnected pieces, $C_3^{\text{disc}}$ from Eq.~\eqref{eq:C3_disc_ens} enters.

In the limit of large source-insertion and insertion-sink separations, excited states
decay exponentially and the ratio plateaus to the ground-state matrix element:

\begin{equation}\label{eq:plateau}
  R_{\mu\nu}^{\mathcal{O}}(q, \tau, t_{\text{sep}}, t_f)
    \xrightarrow[\tau \gg 0,\; t_{\text{sep}}-\tau \gg 0]{}
    \frac{\bra{\pi(p_f)} T_{\mu\nu}^{\mathcal{O}}(t_f) \ket{\pi(p_i)}}
         {2\sqrt{E_\pi(p_i)\,E_\pi(p_f)}}.
\end{equation}

With $p_i = 0$ and $E_\pi(0) = m_\pi$, this simplifies to

\begin{equation}\label{eq:plateau_rest}
  R_{\mu\nu}^{\mathcal{O}} \;\xrightarrow{\text{plateau}}\;
    \frac{\bra{\pi(p_f)} T_{\mu\nu}^{\mathcal{O}}(t_f) \ket{\pi(0)}}
         {2\sqrt{m_\pi\,E_\pi(p_f)}}.
\end{equation}

\textbf{Plateau fitting procedure:}
\begin{enumerate}[nosep]
  \item For each $t_{\text{sep}}$, plot $R(\tau)$ as a function of $\tau$.
  \item Identify the plateau region where $R(\tau)$ is flat within errors.
    Typical choices: $\tau \in [2,\, t_{\text{sep}}-2]$ or $\tau \in [t_{\text{sep}}/3,\, 2t_{\text{sep}}/3]$.
  \item Fit $R(\tau)$ to a constant (or constant $+$ single-exponential correction)
    in the plateau window.
  \item The fit result is the bare matrix element
    $M_{\mu\nu}^{\mathcal{O},\,\text{bare}}(q, t_f, t_{\text{sep}})$.
  \item As a cross-check, verify that the extracted $M$ is independent of $t_{\text{sep}}$
    (within errors) for sufficiently large $t_{\text{sep}}$.
\end{enumerate}

Alternatively, the \textbf{summation method} integrates over $\tau$:

\begin{equation}\label{eq:summation}
  S_{\mu\nu}^{\mathcal{O}}(q, t_{\text{sep}}, t_f)
    = \sum_{\tau = \tau_{\text{min}}}^{t_{\text{sep}} - \tau_{\text{min}}}
      R_{\mu\nu}^{\mathcal{O}}(q, \tau, t_{\text{sep}}, t_f),
\end{equation}

whose slope against $t_{\text{sep}}$ gives the matrix element with reduced excited-state
contamination at the cost of larger statistical errors.

% ------------------------------------------------------------------------------
\subsubsection{Step 7 -- Assemble, renormalize, and decompose}
\label{sec:analysis_step7}

\paragraph{Total bare matrix element.}
After ground-state extraction, the three pieces are summed at each flow time $t_f$:

\begin{equation}\label{eq:M_total_bare}
  M_{\mu\nu}^{\text{bare}}(q, t_f)
    = M_{\mu\nu}^{\text{conn}}(q, t_f)
    + M_{\mu\nu}^{\text{disc},\,q}(q, t_f)
    + M_{\mu\nu}^{\text{disc},\,g}(q, t_f).
\end{equation}

All three pieces share the same lattice kinematics and flow schedule, so they can be
summed directly.

\paragraph{Renormalization and flow-time dependence.}
The bare flowed matrix elements are converted to the $\overline{\text{MS}}$ scheme at scale $\mu$
using perturbative matching coefficients \cite{Suzuki:2013, Makino:2016}:

\begin{equation}\label{eq:renorm_full}
  M_{\mu\nu}^{\overline{\text{MS}}}(q, \mu)
    = \lim_{t_f \to 0}\,
      \Big[ c_q(t_f, \mu)\; M_{\mu\nu}^{q,\,\text{bare}}(q, t_f)
         + c_g(t_f, \mu)\; M_{\mu\nu}^{g,\,\text{bare}}(q, t_f)
         + \delta_{\mu\nu}\, c_{\text{trace}}(t_f, \mu)\,
           \sum_{\rho,\sigma} M_{\rho\sigma}^{\text{bare}}(q, t_f)
      \Big].
\end{equation}

Here $M_{\mu\nu}^{q,\,\text{bare}} = M_{\mu\nu}^{\text{conn}} + M_{\mu\nu}^{\text{disc},\,q}$ is the total quark contribution,
and $M_{\mu\nu}^{g,\,\text{bare}} = M_{\mu\nu}^{\text{disc},\,g}$ is the total gluon contribution.
The trace term accounts for the trace anomaly; $c_{\text{trace}}$ is determined by the
renormalization condition $\sum_\mu T_{\mu\mu}^{\text{ren}} = (\beta/2g) F_{\rho\sigma}^a F_{\rho\sigma}^a
+ (1+\gamma_m) m \bar{\psi}\psi$.

In practice, the $t_f \to 0$ limit is implemented by:
\begin{enumerate}[nosep]
  \item Extracting the matrix element at each flow time $t_f = \code{step} \times \varepsilon$.
  \item Applying the matching coefficients $c_q(t_f, \mu)$, $c_g(t_f, \mu)$ at each $t_f$.
  \item Performing a linear (or constant) extrapolation in $t_f$ to $t_f = 0$,
    using only the small-$t_f$ window where perturbation theory is reliable
    (typically $\sqrt{8t_f} \lesssim 0.3$\,fm).
\end{enumerate}

\paragraph{Gravitational form factor decomposition.}
For a spin-0 pion, the Lorentz decomposition of the renormalized EMT matrix element is
\cite{DelDebbio:2021}:

\begin{equation}\label{eq:GFF_decomp}
  \bra{\pi(p_f)} T_{\mu\nu}^{\overline{\text{MS}}}(\mu) \ket{\pi(p_i)}
    = 2\,P_\mu P_\nu\, A(q^2)
    + \tfrac{1}{2}\big(q_\mu q_\nu - \delta_{\mu\nu} q^2\big)\, D(q^2)
    + 2\,m_\pi^2\,\delta_{\mu\nu}\, \bar{c}(q^2),
\end{equation}

where $P = (p_i + p_f)/2$ and $q = p_f - p_i$.
The three gravitational form factors are:
\begin{itemize}[nosep]
  \item $A(q^2)$ -- the mass form factor ($A(0) = 1$ for a stable hadron),
  \item $D(q^2)$ -- the $D$-term encoding pressure and shear distributions,
  \item $\bar{c}(q^2)$ -- the trace-anomaly form factor ($\bar{c}(0)$ gives the gluon
    contribution to the pion mass).
\end{itemize}

Given $M_{\mu\nu}(q) \equiv \bra{\pi(p_f)} T_{\mu\nu} \ket{\pi(p_i)}$ from the lattice,
the form factors are extracted by solving the $4 \times 4$ linear system at each $q^2$.
With $p_i = 0$, the components simplify:

\begin{align}
  M_{44}(q) &= 2E_\pi^2(q)\, A(q^2)
              + \tfrac{1}{2}\big(q_4^2 + \vec{q}\,^2\big)\, D(q^2)
              + 2 m_\pi^2\, \bar{c}(q^2), \label{eq:M44} \\
  M_{ii}(q) &= 2q_i^2\, A(q^2)
              + \tfrac{1}{2}\big(q_i^2 - \vec{q}\,^2\big)\, D(q^2)
              + 2 m_\pi^2\, \bar{c}(q^2) \quad (i = 1,2,3), \label{eq:Mii} \\
  M_{4i}(q) &= 2E_\pi(q)\,q_i\, A(q^2)
              + \tfrac{1}{2}\,q_4\,q_i\, D(q^2), \label{eq:M4i} \\
  M_{ij}(q) &= 2q_i q_j\, A(q^2)
              + \tfrac{1}{2}\,q_i q_j\, D(q^2) \quad (i \neq j). \label{eq:Mij}
\end{align}

In practice, $A(q^2)$, $D(q^2)$, and $\bar{c}(q^2)$ are overdetermined (10 independent
$M_{\mu\nu}$ components for 3 unknowns at each $q^2$) and are extracted by a least-squares
fit to Eqs.~\eqref{eq:M44}--\eqref{eq:Mij}.

\paragraph{Summary of the full analysis chain.}
\begin{center}
\fbox{
\begin{minipage}{0.92\textwidth}
\vspace{4pt}
\begin{enumerate}[nosep, leftmargin=*]
  \item \textbf{Read} HDF5 files for all $N_{\text{cfg}}$ configurations.
  \item \textbf{Match} momenta: $p_f = q$, locate indices in \code{momentum\_list} and
    \code{qext}.
  \item \textbf{Form} $C_3^{\text{disc},\,q}$ and $C_3^{\text{disc},\,g}$ per
    configuration via Eq.~\eqref{eq:C3_disc_per_cfg}.
  \item \textbf{Jackknife} over configurations; compute $\langle C_2 \rangle$,
    $\langle L \rangle$, $\langle C_2 L \rangle$; subtract vacuum via
    Eq.~\eqref{eq:C3_disc_ens}.
  \item \textbf{Ratio} $R = C_3 / \langle C_2 \rangle$; plateau-fit in $\tau$
    for each $(q, t_f, \mu, \nu, t_{\text{sep}})$.
  \item \textbf{Sum} connected + disconnected quark + gluon $\to$
    $M_{\mu\nu}^{\text{bare}}(q, t_f)$.
  \item \textbf{Renormalize} with $c_q(t_f,\mu)$, $c_g(t_f,\mu)$; extrapolate
    $t_f \to 0$.
  \item \textbf{Solve} for $A(q^2)$, $D(q^2)$, $\bar{c}(q^2)$ via
    Eqs.~\eqref{eq:M44}--\eqref{eq:Mij}.
\end{enumerate}
\vspace{2pt}
\end{minipage}
}
\end{center}

\subsection{Future Variance Reduction Strategies}
\label{sec:variance_reduction}

The current implementation uses plain $\znum$ noise for the stochastic quark 1pt estimator. Two natural upgrades are identified for future production:

\begin{enumerate}[nosep]
  \item \textbf{Hierarchical probing} \cite{Stathopoulos:2013} (arXiv:1302.4018):
  The EMT 1pt trace estimator is $\Tr(\Gamma D^{-1})$ with local gamma and derivative insertions. Hierarchical probing replaces or supplements random noise vectors with distance-ordered Hadamard/coloring vectors on the 4D torus. This cancels near-diagonal contributions of $D^{-1}$ more deterministically, reducing stochastic variance at fixed inversion count. The method is particularly effective for operators with local support like $T_{\mu\nu}$.

  \item \textbf{Frequency splitting / propagator decomposition} (arXiv:2605.00643):
  The flowed EMT loop can be decomposed into low/infrared and high/ultraviolet components using frequency-filtered estimator pieces. This allows expensive noise averaging to be focused on the component with the largest residual variance. The strategy should be designed together with the gradient-flow radius since flow already smooths ultraviolet modes.
\end{enumerate}

When implemented, these upgrades should preserve the existing HDF5 schema or add explicit estimator metadata so that disconnected diagram analysis can combine old and new 1pt data safely.

\subsection{Summary of All HDF5 Output Files}
\label{sec:hdf5_summary}

\begin{table}[h]
\centering
\caption{Summary of HDF5 output files produced by the pion EMT workflow.}
\label{tab:hdf5}
\begin{tabular}{@{}llll@{}}
\toprule
\textbf{File type} & \textbf{Tag subdirectory} & \textbf{Main datasets} & \textbf{Shape} \\
\midrule
Meson 2pt     & \code{EMT2pt/}  & \code{C2}          & \code{[16, Np, Nt]} \\
Quark 3pt     & \code{EMT3pt/}  & \code{C2}          & \code{[Nt]} \\
              &                  & \code{C3\_chi}      & \code{[Nts, Ns+1, Nq, Nt]} \\
              &                  & \code{C3\_Tmunu}    & \code{[Nts, Ns+1, Nq, 4, 4, Nt]} \\
Quark 1pt     & \code{EMTc/}    & \code{raw/derivative\_bilinear\_pervec} & \code{[Nv,16,4,Nq,Ns+1,Nt]} \\
              &                  & \code{avg/Tmunu/Tij}    & \code{[Nq, Ns+1, Nt]} \\
Gluon 1pt       & \code{EMTg/<lat>.EMTg.<cfg>.<ama>.<sm>.h5} & \code{Tmunu/Tij} & \code{[Nq, Ns+1, Nt]} \\
\bottomrule
\end{tabular}
\end{table}

Notation: \code{Np} = number of momenta, \code{Nq} = number of momentum transfers, \code{Nts} = number of $t_{\text{sep}}$ values, \code{Ns} = number of flow steps, \code{Nv} = number of noise vectors, \code{Nt} = temporal lattice extent.

% ==============================================================================
\section{Technical Implementation Details}
% ==============================================================================

\subsection{Parameter Dictionary}
\label{sec:parameters}

The \code{QuarkEMT} class is initialized with a parameter dictionary containing:

\begin{center}
\begin{tabular}{@{}ll@{}}
\toprule
\textbf{Key} & \textbf{Description} \\
\midrule
\code{qext}       & list of momentum transfers $[q_x, q_y, q_z, q_t]$ \\
\code{pf}         & fixed sink momentum $[p_x, p_y, p_z, p_t]$ \\
\code{p\_2pt}     & list of momenta for the 2pt function \\
\code{CG\_GaussSmear} & whether to apply Gaussian smearing \\
\code{pos\_boost} & momentum boost for forward source/sink \\
\code{neg\_boost} & momentum boost for backward source/sink \\
\code{width}      & Gaussian smearing width \\
\code{flow\_type} & \code{"wilson"} or \code{"symanzik"} \\
\code{flow\_epsilon} & flow step size $\varepsilon$ \\
\code{flow\_steps}   & number of flow steps \\
\bottomrule
\end{tabular}
\end{center}

\subsection{Inverter Parameters}
\label{sec:inv_par}

The Dirac operator uses the QUDA multigrid solver with parameters:
\begin{center}
\code{invPara = [mass, csw, tol, maxiter]}
\end{center}
Typical values: \code{mass = 0.236}, \code{csw = 1.0372}, \code{tol = 1e-10}, \code{maxiter = 300}. The multigrid setup uses a $[[8, 8, 4, 4]]$ block size configuration.

\subsection{Random Number Parameters (Quark 1pt)}
\label{sec:rand_par}

\code{randPara = [n\_vec, n\_zn, randseed]} where \code{n\_vec} is the
number of base noise vectors, \code{n\_zn} is the $\znum$ order (default 4),
and \code{randseed} is the counter stream salt (default 0).  The gauge
configuration number is a separate counter key.  No rank-local backend RNG is
used.

\subsection{Typical Usage Pattern}
\label{sec:usage}

The application scripts follow a common pattern:

\begin{enumerate}[nosep]
  \item \textbf{Initialize QUDA}: \code{init(mpi\_geometry, enable\_mps=True)}
  \item \textbf{Read gauge field}: \code{io.readNERSCGauge(gauge\_path)}
  \item \textbf{Apply HYP smearing}: \code{gauge.hypSmear(1, 0.75, 0.6, 0.3, 4)}
  \item \textbf{Set boundary condition}: \code{gauge.latt\_info.t\_boundary = -1} (anti-periodic in time)
  \item \textbf{Instantiate measurement class}: \code{quark\_emt = QuarkEMT(parameters)}
  \item \textbf{Run measurement}: \code{quark\_emt.connected\_3pt(...)}
  \item \textbf{Results saved to HDF5} automatically.
\end{enumerate}


% ==============================================================================
\section{Cross-Checks and Consistency}
% ==============================================================================

\begin{itemize}
  \item \textbf{Convention B + meson\_sign = +1}: All meson 2pt and 3pt contractions are internally consistent with this choice. Zero-momentum regression tests match previous baselines to roundoff precision.
  \item \textbf{CHI diagnostic}: $C_3^\chi$ with $\mathcal{O}_g = \mathbf{1}$ must be numerically consistent with the meson 2pt function (up to a sink momentum phase factor). This is a strong check of the sequential source construction.
  \item \textbf{Symmetrization}: Both connected and disconnected EMT tensors are symmetrized $\mu \leftrightarrow \nu$ after the full contraction. The antisymmetric part should be zero within statistics.
  \item \textbf{Flow schedule alignment}: The quark and gluon flow schedules are aligned (measure-before-flow, first interval subdivided). Output index \code{step} corresponds to flow time $t = \text{step} \times \varepsilon$ in all three file types.
  \item \textbf{Volume normalization}: The connected 3pt output is NOT volume-normalized (it is an extensive correlator). The 1pt outputs ARE divided by $V_3$ (intensive densities). This distinction is critical for the analysis stage when forming $\langle C_2 L \rangle$.
\end{itemize}


% ==============================================================================
\begin{thebibliography}{99}
% ==============================================================================

\bibitem{Luscher:2010iy}
M.\ L\"uscher,
``Properties and uses of the Wilson flow in lattice QCD,''
JHEP \textbf{08} (2010), 071.
[arXiv:1006.4518 [hep-lat]].

\bibitem{Luscher:2013vga}
M.\ L\"uscher and P.\ Weisz,
``Perturbative analysis of the gradient flow in non-abelian gauge theories,''
JHEP \textbf{02} (2011), 051.
[arXiv:1101.0963 [hep-th]].

\bibitem{Makino:2016}
H.\ Makino and H.\ Suzuki,
``Lattice energy-momentum tensor from the Yang-Mills gradient flow
-- inclusion of fermion fields,''
PTEP \textbf{2014} (2014), 063B02.
[arXiv:1404.2758 [hep-lat]].

\bibitem{Stathopoulos:2013}
A.\ Stathopoulos, J.\ Laeuchli, and K.\ Orginos,
``Hierarchical probing for estimating the trace of the matrix inverse on
toroidal lattices,''
[arXiv:1302.4018 [hep-lat]].

\bibitem{DelDebbio:2021}
L.\ Del Debbio, A.\ Patella, and C.\ Pica,
``Energy-momentum tensor on the lattice: recent developments,''
[arXiv:2105.08278 [hep-lat]].

\bibitem{Harris:2020}
T.\ Harris, G.\ von Hippel, P.\ Junnarkar, H.\ B.\ Meyer, K.\ Ottnad,
J.\ Wilhelm, H.\ Wittig, and L.\ Wrang,
``Nucleon isovector tensor form factors from the lattice,''
Phys.\ Rev.\ D \textbf{100} (2019), 034513.
[arXiv:1905.01291 [hep-lat]].

\bibitem{Suzuki:2013}
H.\ Suzuki,
``Energy-momentum tensor from the Yang-Mills gradient flow,''
PTEP \textbf{2013} (2013), 083B03.
[arXiv:1304.0533 [hep-lat]].

\end{thebibliography}

\end{document}
